\section{The remaining Gaussian reflection-component cases}
\label{sec:gaussian-reflection-reduction}

Throughout, the original finite-group representation is faithful,
absolutely irreducible over \(K=\mathbb Q(i)\), and has exact
character field \(K\). We study the actual primitive-factor
polarization, and a sixfold of signature \((3,3)\).

\subsection*{D-GAUSS-01. The Gaussian \(F_4\) normalizer and \(D_4\) triality}

Use the standard rational \(F_4\) roots and standard form on \(K^4\).
Let \(W=W(F_4)\), and put
\[
 \Phi=\begin{pmatrix}0&-1&-1&0\\1&0&0&1\\1&0&0&-1\\0&-1&1&0\end{pmatrix},
 \qquad z=\frac{1+i}{2}\Phi,\qquad X=\langle W,z\rangle.
\]
The matrix \(\Phi/\sqrt2\) interchanges the long and short root
lines. This was verified on the complete 48-root system in D-GF4-01.
Moreover \(\Phi^2=-2I\), so \(z^2=-iI\), and \(z\) normalizes
\(W\). Consequently
\[
       X=W\sqcup iW\sqcup zW\sqcup izW,\qquad |X|=4608.
\]

Every finite unitary normalizer of \(W\) over \(K\) is contained
in \(X\). Here are the details of the scalar issue. The rational
intertwiner argument of D-FIELD-02 writes such an element as \(g=cA\)
with \(A\) rational and \(c\in K^*\). Finite order implies
\(c^N\in\mathbb Q\), hence \(c/\bar c\in\mu_4\).
Therefore \(c\) is a rational multiple of one of
\(1,i,1+i,1-i\). The rational similarity \(A\) either preserves
the two root-length classes or interchanges them. After rational
rescaling it is respectively a Weyl element or a Weyl element times
\(\Phi\), by the chamber argument and the absence of a nontrivial
automorphism of the oriented \(F_4\) diagram. In the first case
unitarity forces the remaining scalar to be in \(\mu_4\). In the
second case it forces \(2N(\lambda)=1\); the listed four possible
rational directions then give
\(\lambda\in\mu_4(1+i)/2\). This proves the containment.

For an irreducible component of type \(D_4\), Cohen's normalizer
Proposition 2.13 gives containment in scalar multiples of \(W(F_4)\),
including the triality automorphisms
\cite[\S2.13, pp.~391--392]{FIELD:Cohen}. Its rational Weyl matrices
and the \(K\)-model force the finite scalar to lie in \(\mu_4\).
Thus every such stabilizer image is contained in \(\mu_4W\subset X\).

An exhaustive exact enumeration from the four standard simple \(F_4\)
reflections and \(z\) gives 58 conjugacy classes in \(X\).
The simple roots used are
\[
 e_2-e_3,\quad e_3-e_4,\quad e_4,\quad
                   (e_1-e_2-e_3-e_4)/2.
\]
The integer program stores twice every Gaussian matrix, verifies every
product's divisibility, enumerates the full generated group, verifies
unitarity, and partitions it by conjugation. The rational checker then
computes every moving-space Gram determinant and local invariant as in
D-E67-01, with exact cyclotomic Fourier spectra. There are just three
nonnorm classes:
\[
\begin{array}{c|c|c|c|c|c}
 \text{certificate index}&\text{size}&\text{order}&k_g&q_g&
                       \text{age}(g)-k_g/2\\\hline
11&96&12&4&3&1/2\\
31&96&12&4&3&-1/2\\
37&96&6&4&3&0
\end{array}
\]
In particular every element of moving rank at most three has norm local
invariant. All matrices, classes and invariants are retained in
\texttt{F4\_gaussian\_classes.json} and
\texttt{F4\_gaussian\_local\_invariants.json}, with the short
programs \texttt{enumerate\_f4\_gaussian.cpp} and
\texttt{check\_f4\_gaussian.py}.

A genus-zero sixfold source in this rank-four representation has
\(D=14\). A nonsplit balanced determinant requires a nonnorm local
factor by D-DET-01, hence an inertia of moving rank four. Its branch
count therefore satisfies \(14\ge4+(r-1)\), or \(r\le11\).
Its period image has dimension at most eight. For a base of genus
one, \(D=6\) bounds the dimension by six; higher genera are impossible.
This argument applies to every subgroup of \(X\), since it uses the
actual matrix local invariants and does not require that the subgroup
contain all classes. In particular it excludes dominant nonsplit
sources from the \(D_4\) and \(F_4\) reflection components, including
both triality and the Gaussian length-interchange extension.

\subsection*{D-GAUSS-02. An all-rank obstruction for the standard \(A_m\) component}

Let \(m\ge3\), \(n=m+1\), and let \(W(A_m)=S_n\) act on its
standard sum-zero subspace. The finite Gaussian unitary normalizer is
contained in \(\mu_4W(A_m)\): all roots have the same length, and
the chamber argument identifies root-system automorphisms with Weyl
elements and the diagram involution \(-w_0\). The rational-intertwiner
and scalar argument used in D-E67-02 now applies. A stabilizer image
containing \(W(A_m)\) and having nonrational character therefore
contains \(iI\) and is the whole group
\[
                J=\mu_4\times S_n,\qquad
                \rho=\chi_i\otimes\operatorname{Std}_n.
\]
The product is direct because no nonidentity scalar occurs in the
standard \(S_n\)-representation when \(n\ge3\).
This standard representation is not being assumed monomial.

Consider first a genus-zero generating branch tuple. Write its elements
as \(\lambda\pi\), with \(\lambda\in\mu_4\), \(\pi\in S_n\).
If \(\lambda\) has order four, the fixed dimension of
\(\lambda\pi\) on the permutation representation is the number
of cycles of \(\pi\) whose lengths are divisible by four. The
constant line has eigenvalue \(\lambda\ne1\), so removing it does
not change the fixed dimension. Thus
\[
      k(\lambda\pi)\ge m-\lfloor n/4\rfloor.
\]
Likewise, for \(\lambda=-1\),
\[
      k(-\pi)\ge m-\lfloor n/2\rfloor.
\]
Let \(t\) be the number of order-four scalar entries and \(u\)
the number of scalar-minus-one entries. Generation and the scalar
product-one relation imply that \(t\ge2\) is even.
Every branch is nonidentity; a rank-nine period map would require
\(r\ge12\), and hence the total moving excess would satisfy
\[
    \sum_j(k_j-1)=2m+6-r\le2m-6=2n-8.
\]
If \(t\ge2\) and \(u\ge1\), its lower bound from these three
entries already exceeds the right side: the difference is
\[
 3(n-2)-2\lfloor n/4\rfloor-\lfloor n/2\rfloor-(2n-8)
      =n+2-2\lfloor n/4\rfloor-\lfloor n/2\rfloor\ge2.
\]
If \(t\ge4\), the four order-four entries alone exceed the bound,
with difference
\[
              4(n-2-\lfloor n/4\rfloor)-(2n-8)
                  =2n-4\lfloor n/4\rfloor>0.
\]
Consequently a possible rank-nine source has
\[
                  t=2,\qquad u=0.
\]
The two nontrivial scalars must be \(i\) and \(-i\).

Let \(C\) be the connected \(J\)-cover. Its cyclic quotient
\(T=C/S_n\to\mathbb P^1\) has exactly two branch points, both
fully ramified of order four. Riemann--Hurwitz gives \(g(T)=0\).
Put
\[
 H=\mu_4\times S_{n-1},\qquad
 X=C/S_{n-1},\qquad Y=C/H.
\]
Then \(X\to Y\) is a connected cyclic degree-four cover. The
representation identity
\[
 \operatorname{Ind}_H^J(\chi_i)
       =\chi_i\otimes\operatorname{Perm}_n
       =\chi_i\oplus\rho
\]
and Frobenius reciprocity identify its faithful-character multiplicity
space with the direct sum of those for \(\chi_i\) and \(\rho\).
The first is zero, because it lies in \(H^1(T)\). Applying the same
argument to conjugate characters gives a rational Hodge isogeny
\[
                        P_4(X/Y)\sim B_\rho.
\]
It is induced by fixed group-algebra correspondences in the family.

This identification also controls the actual polarization. The
Rosati-fixed rational projector for the pair \(\chi_i,\bar\chi_i\)
on \(H\) has one primitive line in the \(\rho\)-block and a line
in the scalar-character block. The latter acts as zero on \(H^1(C)\).
Its nonzero abelian image is therefore a primitive \(\rho\)-factor,
and is the pullback of the cyclic factor from \(X\). Pullback
multiplies the Jacobian polarization by \(|S_{n-1}|\). Any other
Rosati-fixed primitive line changes the multiplicity Hermitian form by
a positive rational scalar, as in D-MON-03. Hence splitness of this
cyclic factor is splitness of the specified primitive polarization.

If \(h=g(Y)\) and \(R\) is the cyclic branch count, the faithful
parabolic dimension is now exactly six:
\[
                              6=2h-2+R.
\]
The actual cyclic variation has moduli dimension at most
\(5,6,7,8,9\) for \(h=0,1,2,3,4\). Its isogeny to the
sixfold variation preserves the period differential's kernel. Thus
rank nine forces \(h=4,R=0\). Looijenga's equivariant hyperbolic
basis, already verified and applied in D-MON-02--03, makes the actual
cyclic-factor polarization split. This excludes the desired nonsplit
source for every \(m\), without invoking monomiality of the standard
representation.

Higher-genus original bases cannot give rank nine. At genus one,
\(D=6\), so their dimension is at most six. At genus two,
\(D=6-2m\), leaving only \(m=3,D=0\) among the ranks under
consideration, with base-moduli dimension three. Larger genus is
impossible. Ranks one and two are already excluded respectively by
D-MON and the rank-two Hurwitz-dimension bound in D-FIELD-03.
Induction from a reflection-component stabilizer preserves the actual
variation and polarization as in that same statement. Therefore the
obstruction also applies to arbitrary original finite groups with an
\(A_m\) reflection component.

\subsection*{D-GAUSS-03. Exact scope of the remaining reflection-containing gate}

Let \(R\) be the normal closure of the reflections in the original
irreducible Gaussian group. As in D-FIELD-02, its nonorthogonality
components are defined over \(K\), the original group permutes them
transitively, and the original representation is induced from the
stabilizer of one irreducible reflection component.

Cohen's classification and Malle's character-field tables show that
the nonreal primitive reflection components over \(K\) are \(G_8\)
in rank two and \(G_{29},G_{31}\) in rank four
\cite[\S3.6 and Theorem 5.12]{FIELD:Cohen}
\cite[Tables 8.1--8.2, pp.~674--675]{FIELD:Malle}.
The omissions \(G_4,G_{25}\) in Malle's table have order-three
reflections and are not Gaussian; the omitted \(G_{28},G_{35}\)
are the rational Weyl groups \(F_4,E_6\).
The primitive real possibilities are rational Weyl groups, since the
noncrystallographic real character fields are not contained in \(K\).

An imprimitive component has a unique system of imprimitivity except
for the four cases in Cohen's Lemma 2.7
\cite[pp.~389--390]{FIELD:Cohen}. The rank-two exceptions are already
excluded by their dimension bound; \(G(3,3,3)\) has non-Gaussian
character field; and \(G(2,2,4)=W(D_4)\) is covered by D-GAUSS-01.
In every other imprimitive case its normalizer preserves that unique
system, so its representation is monomial. Transitivity of induction
then makes the original representation monomial, and D-MON excludes
dominant nonsplit sources even if the inducing character has a larger
cyclotomic field. This covers the usual \(B_m,D_m\) cases as well.

For \(G_{29},G_{31}\), the corrected normalizer Proposition 5.14
of Cohen puts the stabilizer image in scalar multiples of that group
\cite{FIELD:Cohen,FIELD:CohenErr}. Each already contains \(\mu_4\),
as is visible in the central classes of the exact enumeration in
D-G31-01. The scalar lies in \(K\), and finite order forces it
into \(\mu_4\). Thus the stabilizer image is the same natural
group, and D-G31-02 applies.

The remaining rational Weyl components are accounted for by
D-GAUSS-01 (\(D_4,F_4\)), D-GAUSS-02 (all \(A_m\)),
D-E67-02 (\(E_6,E_7\)), and the rank-two dimension bound.
Accordingly, among reflection-containing Gaussian sources in the stated
group-algebra setup, the sole remaining possible component for a
dominant nonsplit sixfold is
\[
                              \mu_4W(E_8).
\]
This is a reduction to a candidate component, not a dominance theorem
or an algebraic-cycle supplier. Reflection-free representations remain
outside this classification and must be considered separately.
